
\subsection{Flow-based generative models in physics applications}

Flow-based generative models are increasingly used in physics for probabilistic forecasting, learning to sample from the conditional distribution of the next state given the current one. Some approaches combine deterministic neural operators with noise-based techniques: \cite{lippe_pde-refiner_2023} add/remove noise around neural-operator predictions, while \cite{gao_generative_2024} time-step in latent space followed by physics-informed diffusion upsampling. Others target direct conditional sampling: \cite{kohl_benchmarking_2024} train denoising diffusion models, \cite{fotiadis_stochastic_2024} use flow matching, and \cite{chen_probabilistic_2024} use stochastic interpolants, which \cite{mucke_physics-aware_2025} extended with energy consistency in training.

\subsection{Posterior sampling with flow-based generative models}
Generative models are increasingly used for inverse problems \cite{daras_survey_2024}. Several training-free methods address linear inverse problems: \cite{zhang_flow_2025} matches corrupted trajectories along flow priors, \cite{ahamed_dawn-si_2024} (DAWN-SI) uses data-aware stochastic interpolation, and \cite{negrel_multitask_2025} explores multitask learning. A second line conditions a deterministic generator on measurements: the decomposition $\nabla\log p_\tau(\bx\mid\obs)=\nabla\log p_\tau(\bx)+\nabla\log p_\tau(\obs\mid\bx)$ of \cite{song_score-based_2021} underlies score-based conditional sampling and equally defines a conditional probability-flow ODE. For flow matching, OT-ODE \cite{pokle_training-free_2024} and FIG \cite{yan_fig_2024} guide the probability-flow ODE -- FIG with a measurement interpolant coinciding with ours \eqref{eq:interpolated_observation} -- while \cite{zhang_flow_2025} solves a MAP problem along the flow and \cite{martin_pnp-flow_2025} folds the prior into a plug-and-play denoiser. The Gaussian likelihood surrogate these methods (and ours, Section~\ref{subsec:obs_interpolation}) rely on originates in diffusion guidance: DPS \cite{chung_diffusion_2023} uses the denoiser mean with uninflated covariance, $\Pi$GDM \cite{song_pseudoinverse_2023} inflates it. The likelihood-score coefficient in these ODE methods is fixed by the velocity--score conversion \cite{albergo_stochastic_2023, ma_sit_2024}, exactly the duality coefficient $\vscoef_\tau$ in our guided-ODE weight. Our guided ODE generalises these to autoregressive DA, as the zero-diffusion member of the family containing the SDE samplers.

\subsection{Data assimilation with generative models}
Generative DA seeks to overcome the limitations of ensemble Kalman and particle filters. Early efforts include score-based DA \cite{rozet_score-based_2023}; for high-dimensional systems \cite{bao_ensemble_2024} introduced an ensemble score filter, extended to latent space for sparse observations by \cite{si_latent-ensf_2024}, with multimodal settings addressed by \cite{qu_deep_2024}. Most relevant is FlowDAS \cite{chen_flowdas_2025}, a stochastic-interpolant DA framework with a Monte-Carlo likelihood estimate.

% \subsection{From model-specific to model-agnostic assimilation}
% The methods above are typically tied to a particular generator: score-based and stochastic-interpolant methods exploit a diffusion term, whereas flow matching methods \cite{pokle_training-free_2024, zhang_flow_2025, yan_fig_2024} correct a deterministic ODE. We unify these by viewing both as affine Gaussian probability paths \cite{albergo_stochastic_2023, lipman_flow_2023} and conditioning that path on the observation. Two known equivalences make this possible: velocity and score determine each other in closed form \cite{albergo_stochastic_2023, ma_sit_2024}, and the probability-flow ODE is one member of a family of equal-marginal SDEs with freely tunable diffusion \cite{song_score-based_2021}. Applied to the conditioned path, these yield the conditional score and velocity from a single likelihood-score term, and the equal-marginal family produces the SDE samplers and the guided ODE together, rather than a separate guidance scheme for each. Relative to FlowDAS \cite{chen_flowdas_2025}, which commits to a stochastic-interpolant SDE with a Monte-Carlo likelihood estimate, we add the lifted DM-SDE and the deterministic guided ODE, and replace the sampling-based likelihood with a closed-form interpolated-observation likelihood with inflated covariance.
